\section*{الفصل \ref{ch:b2:geological-hazards} --- الأخطار الجيولوجية ومخاطرها}

\begin{solution}{exo:b2:geological-hazards:1}
(أ) $1 - e^{-30/100} = 0.259$. (ب) $e^{-100/100} = e^{-1} = 0.368$: فأكثر من قرن واحد من كلّ ثلاثة قرون
يمرّ من دون فيضان كهذا. (ج) $T = -50/\ln 0.95 = 975$ سنة.
\end{solution}

\begin{solution}{exo:b2:geological-hazards:2}
(أ) البلدة A: $(1/200) \times 5000 \times 0.02 = 0.5$ مبنى في السنة؛ والبلدة B: $(1/50) \times 800
\times 0.10 = 1.6$. (ب) البلدة B، مع أنّ عدد مبانيها المعرَّضة أقلّ بمقدار 6.25 مرّة. (ج) خطرها أكثر تواترًا
بأربع مرّات وهشاشتها أعلى بخمس مرّات: أي عامل 20 معًا، وهو يرجح على
تعرّضها الأصغر.
\end{solution}

\begin{solution}{exo:b2:geological-hazards:3}
(أ) $\sqrt{9.81 \times 4000} = \qty{198}{m.s^{-1}}$؛ $\sqrt{9.81 \times 200} = \qty{44.3}{m.s^{-1}}$؛
$\sqrt{9.81 \times 10} = \qty{9.90}{m.s^{-1}}$. (ب) $\num{1.5e6}/198.1 = \qty{7570}{s}$، أي \qty{2.10}{h}.
(ج) $800/3.6 = \qty{222.2}{m.s^{-1}}$، و$h = c^2/g = 222.2^2/9.81 = \qty{5030}{m}$.
\end{solution}

\begin{solution}{exo:b2:geological-hazards:4}
(أ) $\tan\ang{33}/\tan\ang{25} = 0.6494/0.4663 = 1.39$. (ب) $F = 1$ عندما $\beta = \varphi' =
\ang{33}$. (ج) مع $c' = 0$، يكون $F = (\cos^2\beta - r_u)\tan\varphi'/(\sin\beta\cos\beta)$، و$\cos^2\beta
- r_u = \cos^2\beta\,(1 - \gamma_w/\gamma)$، ومن ثمّ $F = (1 - \gamma_w/\gamma)\tan\varphi'/\tan\beta$:
فالماء ينصّف الاحتكاك. $(1 - 9.81/20) \times 1.393 = 0.510 \times 1.393 = 0.710$: فينهار
المنحدر.
\end{solution}

\begin{solution}{exo:b2:geological-hazards:5}
(أ) $696/50 = 13.9$ في السنة. (ب) $e^{-30/26.2} = 0.318$، مقابل 0.304 مرصودة: قيمتان متقاربتان. والقيمة
المرصودة أدنى قليلًا لأنّ الهزّات اللاحقة والأزواج الزلزالية المتجمّعة تضيف فواصل
قصيرة. (ج) $1 - e^{-7/26.2} = 0.234$.
\end{solution}

\begin{solution}{exo:b2:geological-hazards:6}
(أ) حدّ القدر هو $-0.105 + 0.229 \times 0.5 = 0.0095$. وعندما $d = 0$ و10 و30
و\qty{100}{km}، يكون $r = 5.57$ و11.45 و30.51 و\qty{100.2}{km}، و$\log_{10}\mathrm{PGA} = -0.571$
و$-0.814$ و$-1.145$ و$-1.547$، ووسيط أقصى التسارع 0.269 و0.153 و0.0715 و\qty{0.0284}{g}.
(ب) $0.153 \times 10^{0.251} = 0.153 \times 1.78 = \qty{0.272}{g}$. (ج) انحراف معياري واحد
على كلّ جانب: من $0.153/1.70 = \qty{0.0900}{g}$ إلى $0.153 \times 1.70 = \qty{0.260}{g}$.
\end{solution}

\begin{solution}{exo:b2:geological-hazards:7}
(أ) بالترتيب: 420 ($T = 13/1 = 13$ سنة)، و305 (6.5)، و268 (4.33)، و251 (3.25). (ب) $305 + (420 -
305)\ln(10/6.5)/\ln(13/6.5) = 305 + 115 \times 0.621 = \qty{376}{m^3.s^{-1}}$. (ج) تُرسم أكبر ذروة
عند 13 سنة؛ وأيّ قيمة عند 100 سنة استقراء بعيد خارج السجلّ، بقانون للقيم القصوى
لا تكاد النقاط الاثنتا عشرة تقيّد اختياره ومعاملاته.
\end{solution}

\begin{solution}{exo:b2:geological-hazards:8}
(أ) $0.6 \times 0.4 \times 0.5 \times 0.3 = 0.036$. (ب) $P(\text{ثوران}) = 0.6 \times 0.4 = 0.24$؛
ومن دون تدفّق في الوادي، $0.24 - 0.036 = 0.204$. (ج) $0.6 \times 0.7 \times 0.5 \times 0.3 =
0.063$.
\end{solution}

\begin{solution}{exo:b2:geological-hazards:9}
في هذا التمرين يكون $\sin\beta\cos\beta = 0.4698$، و$\cos^2\beta = 0.6710$ و$\tan\ang{30}/\tan\ang{35} =
0.8245$. (أ) تعطي المعادلة $F = c'/(\gamma z\sin\beta\cos\beta) + 0.8245 = 1$ القيمة $z_c = 8/(19 \times 0.4698
\times 0.1755) = \qty{5.11}{m}$. (ب) $8/(19 \times 2 \times 0.4698) + 0.8245 = 0.448 + 0.825 =
1.27$. (ج) $r_u = 9.81 \times 0.6710/19 = 0.3465$؛ $F = [8 + 38 \times (0.6710 - 0.3465) \times
0.5774]/(38 \times 0.4698) = 15.12/17.85 = 0.847$: فالطبقة المشبعة تنهار.
\end{solution}

\begin{solution}{exo:b2:geological-hazards:10}
(أ) $\Pr(\text{لا حدث في } [0, s+t] \mid \text{لا حدث في } [0, s]) = e^{-(s+t)/T}/e^{-s/T} = e^{-t/T}$.
(ب) السنوات الأربعمئة الهادئة لا تغيّر شيئًا: $1 - e^{-50/475} = 0.0999$، أي \qty{10}{\percent}. (ج) إنّه
يناقض فرض أنّ المعدّل لا يتوقّف على الزمن المنقضي: فبعد الارتداد يجب أن يتراكم
الإجهاد من جديد، فيكون المعدّل منخفضًا بُعيد الزلزال ثم يزداد لاحقًا. وتضمّ منطقة المصدر
صدوعًا كثيرة في مراحل مختلفة من دوراتها؛ ومجموع متتاليات كثيرة مستقلّة
من الأحداث يسلك سلوكًا قريبًا من عملية بواسون، والمعدّل هو كلّ ما يُعرف عن كلّ منها.
\end{solution}

\begin{solution}{exo:b2:geological-hazards:11}
(أ) $A = 0.5 \times (4000/10)^{1/4} = 0.5 \times 4.47 = \qty{2.24}{m}$. (ب) $\lambda = cT$:
$198.1 \times 1200 = \qty{238}{km}$ في المحيط العميق، و$9.90 \times 1200 = \qty{11.9}{km}$ في
\qty{10}{m} من الماء. (ج) لم تعد السعة صغيرة مقارنة بالعمق (\qty{2.2}{m} في
\qty{10}{m})، فتشتدّ الموجة انحدارًا وقد تنكسر؛ وينعكس جزء من الطاقة أو يتركّز بفعل
شكل قاع البحر والخلجان؛ وعلى اليابسة يتوقّف ارتفاع الصعود على ميل الأرض
وخشونتها.
\end{solution}

\begin{solution}{exo:b2:geological-hazards:12}
(أ) $10/5.8 + 4 = 1.72 + 4 = \qty{5.72}{s}$. (ب) تبلغ موجة S مسافة \qty{50}{km} بعد $50/3.2 =
\qty{15.6}{s}$: أي $15.6 - 5.72 = \qty{9.90}{s}$ من الإنذار؛ وعند \qty{100}{km}، $31.3 - 5.72 =
\qty{25.5}{s}$. (ج) $3.2 \times 5.72 = \qty{18.3}{km}$.
\end{solution}

\begin{solution}{pb:b2:geological-hazards:1}
تُحفَظ اللوغاريتمات بأربعة أرقام عشرية في الخطوات الوسيطة.
\textbf{1.} $10^{4-5} = 0.1$، و$10^{4-6} = 0.01$ و$10^{4-7} = 0.001$ في السنة.
\textbf{2.} 10 و100 و1000 سنة.
\textbf{3.} $10^{4-7.5} = \num{3.16e-4}$ في السنة؛ $0.001 - \num{3.16e-4} = \num{6.84e-4}$ في السنة.
\textbf{4.} $0.1 - 0.01 = 0.09$ و$0.01 - 0.001 = 0.009$ في السنة.
\textbf{5.} $1 - e^{-0.01 \times 50} = 0.393$.
\textbf{6.} $1 - e^{-0.001 \times 50} = 0.0488$.
\textbf{7.} $r = \sqrt{10^2 + 5.57^2} = \qty{11.45}{km}$؛ $\log_{10} r = 1.0588$ و$-0.778 \times
1.0588 = -0.8237$.
\textbf{8.} $\mu = -0.105 + 0.229(M - 6) - 0.8237$: $-1.0432$ و$-0.8142$ و$-0.6424$؛ والأوساط
0.0905 و0.153 و\qty{0.228}{g}.
\textbf{9.} $10^{-0.8142 + 0.230} = \qty{0.260}{g}$.
\textbf{10.} $z = (\log_{10} 0.3 + 0.8142)/0.230 = (-0.5229 + 0.8142)/0.230 = 1.266$، و$1 -
\Phi(1.266) = 0.103$.
\textbf{11.} $-0.105 + 0.229 \times 1.5 - 0.8237 = -0.5852$: أي \qty{0.260}{g}.
\textbf{12.} احتمالات تجاوز \qty{0.2}{g} ($\log_{10} 0.2 = -0.6990$) هي 0.0672
($z = 1.497$)، و0.308 ($z = 0.501$) و0.597 ($z = -0.246$): $\lambda = 0.09 \times 0.0672 + 0.009
\times 0.308 + \num{6.84e-4} \times 0.597 = \num{6.05e-3} + \num{2.77e-3} + \num{4.08e-4} =
\num{9.23e-3}$ في السنة.
\textbf{13.} عند \qty{0.3}{g}: $z = 2.262$ و1.266 و0.520، والاحتمالات 0.0118 و0.103
و0.302؛ $\lambda = \num{1.07e-3} + \num{9.24e-4} + \num{2.06e-4} = \num{2.20e-3}$ في السنة.
\textbf{14.} عند \qty{0.4}{g} ($\log_{10} 0.4 = -0.3979$): $z = 2.805$ و1.810 و1.063،
والاحتمالات 0.00251 و0.0352 و0.144؛ $\lambda = \num{2.26e-4} + \num{3.17e-4} + \num{9.84e-5} =
\num{6.41e-4}$ في السنة.
\textbf{15.} $-\ln 0.9/50 = \num{2.107e-3}$ في السنة، أي فترة عودة قدرها 475 سنة.
\textbf{16.} بين \qty{0.3}{g} (\num{2.20e-3}) و\qty{0.4}{g} (\num{6.41e-4}).
\textbf{17.} الميل هو $\ln(6.41/22.0)/\ln(4/3) = -1.233/0.2877 = -4.29$، و$a = 0.3 \times
(2.107/2.20)^{-1/4.29} = 0.3 \times 1.0101 = \qty{0.303}{g}$.
\textbf{18.} عند \qty{0.3}{g}، تعطي الفئة 5--6 نسبة \qty{48.5}{\percent} من المعدّل، والفئة 6--7
\qty{42.1}{\percent} والفئة 7--7.5 \qty{9.4}{\percent} فقط. وحتى أكبر زلزال وسيطه
\qty{0.260}{g}، أي دون \qty{0.3}{g}: فلا يُبلَغ هذا المستوى إلا في الذيل العلوي للتشتّت،
والزلازل المتوسّطة المتواترة هي التي توفّر معظمه.
\textbf{19.} $4/(19 \times 3 \times 0.4330) + \tan\ang{34}/\tan\ang{30} = 0.162 + 1.168 = 1.33$.
\textbf{20.} $r_u = 9.81 \times 0.75/19 = 0.387$؛ $F = [4 + 57 \times (0.75 - 0.387) \times
0.6745]/(57 \times 0.4330) = 17.95/24.68 = 0.727$: فالمنحدر المشبع ينهار.
\textbf{21.} إذا كان منسوب المياه الجوفية عند كسر $m$ من السمك، فإنّ $r_u = 0.387\,m$. ويكون $F = 1$ عندما
$0.75 - 0.387\,m = (24.68 - 4)/(57 \times 0.6745) = 0.538$: $m = 0.548$، أي منسوب مياه جوفية عند
\qty{55}{\percent} من الطبقة.
\textbf{22.} بقسمة $F = 1$ على $\gamma z$: $0.0702 + (0.75 - 0.4330\,k) \times 0.6745 = 0.4330 +
0.75\,k$، ومن ثمّ $k_c = (0.0702 + 0.5059 - 0.4330)/(0.2921 + 0.75) = 0.1431/1.042 = 0.137$.
\textbf{23.} $\log_{10} 0.137 = -0.8633 = -0.105 + 0.229(M - 6) - 0.8237$ تعطي $M = 6.29$. وتقع زلازل
كهذه بمعدّل $10^{4 - 6.29} - 10^{-3.5} = \num{4.81e-3}$ في السنة، أي نحو مرّة كلّ 208 سنوات،
وأقصى التسارع \omterm{def:b2:geological-hazards:return}{لفترة العودة} 475 سنة في السؤال~17 يفوق ضعف القيمة الحرجة \qty{0.137}{g}: ففي هذا
التحليل ينهار المنحدر الجافّ عند اهتزاز أضعف بكثير من حركة \omterm{def:b2:geological-hazards:return}{فترة العودة} 475 سنة.
\textbf{24.} \omterm{def:b2:geological-hazards:psha}{أقصى تسارع للأرض} له فرصة \qty{10}{\percent} للتجاوز خلال 50
سنة عند البلدة هو \textbf{\qty[text-series-to-math=true]{0.303}{g}}.
\end{solution}
